% Options for packages loaded elsewhere
% Options for packages loaded elsewhere
\PassOptionsToPackage{unicode}{hyperref}
\PassOptionsToPackage{hyphens}{url}
%
\documentclass[
  ignorenonframetext,
]{beamer}
\newif\ifbibliography
\usepackage{pgfpages}
\setbeamertemplate{caption}[numbered]
\setbeamertemplate{caption label separator}{: }
\setbeamercolor{caption name}{fg=normal text.fg}
\beamertemplatenavigationsymbolsempty
% remove section numbering
\setbeamertemplate{part page}{
  \centering
  \begin{beamercolorbox}[sep=16pt,center]{part title}
    \usebeamerfont{part title}\insertpart\par
  \end{beamercolorbox}
}
\setbeamertemplate{section page}{
  \centering
  \begin{beamercolorbox}[sep=12pt,center]{section title}
    \usebeamerfont{section title}\insertsection\par
  \end{beamercolorbox}
}
\setbeamertemplate{subsection page}{
  \centering
  \begin{beamercolorbox}[sep=8pt,center]{subsection title}
    \usebeamerfont{subsection title}\insertsubsection\par
  \end{beamercolorbox}
}
% Prevent slide breaks in the middle of a paragraph
\widowpenalties 1 10000
\raggedbottom
\AtBeginPart{
  \frame{\partpage}
}
\AtBeginSection{
  \ifbibliography
  \else
    \frame{\sectionpage}
  \fi
}
\AtBeginSubsection{
  \frame{\subsectionpage}
}
\usepackage{iftex}
\ifPDFTeX
  \usepackage[T1]{fontenc}
  \usepackage[utf8]{inputenc}
  \usepackage{textcomp} % provide euro and other symbols
\else % if luatex or xetex
  \usepackage{unicode-math} % this also loads fontspec
  \defaultfontfeatures{Scale=MatchLowercase}
  \defaultfontfeatures[\rmfamily]{Ligatures=TeX,Scale=1}
\fi
\usepackage{lmodern}

\usetheme[]{metropolis}
\ifPDFTeX\else
  % xetex/luatex font selection
\fi
% Use upquote if available, for straight quotes in verbatim environments
\IfFileExists{upquote.sty}{\usepackage{upquote}}{}
\IfFileExists{microtype.sty}{% use microtype if available
  \usepackage[]{microtype}
  \UseMicrotypeSet[protrusion]{basicmath} % disable protrusion for tt fonts
}{}
\makeatletter
\@ifundefined{KOMAClassName}{% if non-KOMA class
  \IfFileExists{parskip.sty}{%
    \usepackage{parskip}
  }{% else
    \setlength{\parindent}{0pt}
    \setlength{\parskip}{6pt plus 2pt minus 1pt}}
}{% if KOMA class
  \KOMAoptions{parskip=half}}
\makeatother


\usepackage{longtable,booktabs,array}
\usepackage{calc} % for calculating minipage widths
\usepackage{caption}
% Make caption package work with longtable
\makeatletter
\def\fnum@table{\tablename~\thetable}
\makeatother
\usepackage{graphicx}
\makeatletter
\newsavebox\pandoc@box
\newcommand*\pandocbounded[1]{% scales image to fit in text height/width
  \sbox\pandoc@box{#1}%
  \Gscale@div\@tempa{\textheight}{\dimexpr\ht\pandoc@box+\dp\pandoc@box\relax}%
  \Gscale@div\@tempb{\linewidth}{\wd\pandoc@box}%
  \ifdim\@tempb\p@<\@tempa\p@\let\@tempa\@tempb\fi% select the smaller of both
  \ifdim\@tempa\p@<\p@\scalebox{\@tempa}{\usebox\pandoc@box}%
  \else\usebox{\pandoc@box}%
  \fi%
}
% Set default figure placement to htbp
\def\fps@figure{htbp}
\makeatother





\setlength{\emergencystretch}{3em} % prevent overfull lines

\providecommand{\tightlist}{%
  \setlength{\itemsep}{0pt}\setlength{\parskip}{0pt}}



 


\setbeamerfont{title}{size=\large}
\usepackage{hyperref}
\setbeamertemplate{footline}[frame number]
\usepackage{tikz}
\usetikzlibrary{arrows.meta, positioning, shapes.geometric, fit, backgrounds, calc}
\usepackage{booktabs}
\usepackage{xcolor}
\definecolor{mblue}{RGB}{35,55,100}
\makeatletter
\@ifpackageloaded{caption}{}{\usepackage{caption}}
\AtBeginDocument{%
\ifdefined\contentsname
  \renewcommand*\contentsname{Table of contents}
\else
  \newcommand\contentsname{Table of contents}
\fi
\ifdefined\listfigurename
  \renewcommand*\listfigurename{List of Figures}
\else
  \newcommand\listfigurename{List of Figures}
\fi
\ifdefined\listtablename
  \renewcommand*\listtablename{List of Tables}
\else
  \newcommand\listtablename{List of Tables}
\fi
\ifdefined\figurename
  \renewcommand*\figurename{Figure}
\else
  \newcommand\figurename{Figure}
\fi
\ifdefined\tablename
  \renewcommand*\tablename{Table}
\else
  \newcommand\tablename{Table}
\fi
}
\@ifpackageloaded{float}{}{\usepackage{float}}
\floatstyle{ruled}
\@ifundefined{c@chapter}{\newfloat{codelisting}{h}{lop}}{\newfloat{codelisting}{h}{lop}[chapter]}
\floatname{codelisting}{Listing}
\newcommand*\listoflistings{\listof{codelisting}{List of Listings}}
\makeatother
\makeatletter
\makeatother
\makeatletter
\@ifpackageloaded{caption}{}{\usepackage{caption}}
\@ifpackageloaded{subcaption}{}{\usepackage{subcaption}}
\makeatother

\usepackage{bookmark}
\IfFileExists{xurl.sty}{\usepackage{xurl}}{} % add URL line breaks if available
\urlstyle{same}
\hypersetup{
  pdfauthor={Giorgio Arcara},
  hidelinks,
  pdfcreator={LaTeX via pandoc}}


\author{Giorgio Arcara}
\date{}

\begin{document}


\begin{frame}
\title{Metodologia della Ricerca Clinica\\
\vspace{0.5em}
{\normalsize \textbf{Recap}}}
\author{Giorgio Arcara,\\ Università di Padova \\ IRCCS San Camillo, Venezia}
\titlegraphic{
\vspace*{7cm}
\includegraphics[scale=0.15]{Figures/LogoSanCamilloIRCCS_Unipd_alpha.png}
\hfill \includegraphics[scale=0.15]{Figures/CC_license_3_0.png}
}
\maketitle
\end{frame}

\begin{frame}{1) Test e psicometria}
\phantomsection\label{test-e-psicometria}
Per una pratica psicologica/neuropsicologica evidence based è essenziale
avere \textbf{misurazioni}.

I \textbf{test} sono gli strumenti di misurazione psicologica.
Misurazione e test sono oggetti di studio della \textbf{psicometria}.

Nell'utilizzare nella pratica clinica i test abbiamo un problema
particolare: fare inferenze a livello di \textbf{caso singolo}, che
richiedono attenzioni particolari.
\end{frame}

\begin{frame}{Psicometria applicata}
\phantomsection\label{psicometria-applicata}
Le conoscenze di \textbf{psicometria applicata} entrano in gioco
essenzialmente in due momenti:

\begin{itemize}
\item
  quando scegliamo i test da utilizzare
\item
  quando interpretiamo i risultati dei test
\end{itemize}
\end{frame}

\begin{frame}{Proprietà psicometriche dei test}
\phantomsection\label{proprietuxe0-psicometriche-dei-test}
Alcune considerazioni importanti sono:

\begin{itemize}
\item Considerate \textbf{tutte le proprietà psicometriche di un test}

\begin{itemize}
  \item validità 
  \item affidabilità
  \item dati normativi
  \item cambiamento nel tempo
  \item \ldots
\end{itemize}

\item nella trasformazione di punteggi in percentili (ed eventuale confronto con soglie) ricordate che rispondono alla domanda:

\begin{itemize}
\item Quanto è "strano" osservare questa performance nel mio campione/popolazione normativa?
\end{itemize}

\end{itemize}
\end{frame}

\begin{frame}{Psicometria ed evidenze scientifiche}
\phantomsection\label{psicometria-ed-evidenze-scientifiche}
\includegraphics[width=0.1\linewidth,height=\textheight,keepaspectratio]{Figures/triangle.png}.
In generale ricordatevi che un \textbf{test} \(\nRightarrow\)
\textbf{costrutto}.

\textbf{Dipende dalla validità del test!} (che potrebbe anche cambiare
es. con danni cognitivi)
\end{frame}

\begin{frame}{Studi scientifici}
\phantomsection\label{studi-scientifici}
Per una pratica psicologica/neuropsicologica evidence based è essenziale
basarsi sui risultati di studi scientifici.

La natura delle evidenze può essere \textbf{correlazionale} (tutti gli
studi osservazionali e alcuni interventistici) o \textbf{causale} (trial
clinici randomizzati).
\end{frame}

\begin{frame}{Studi sperimentali e trial}
\phantomsection\label{studi-sperimentali-e-trial}
I \textbf{trial clinici}, anche se sono il metodo principe per efficacia
di un trattamento, \emph{non sono privi di limiti}.

In psicologia clinica e neuropsicologia clinica questo dipende
principalmente da:

\begin{itemize}
\tightlist
\item
  bias nella selezione.
\item
  difficoltà nell'avere un gruppo di controllo appropriato.
\item
  difficoltà nell'avere un \emph{double-blind}.
\end{itemize}
\end{frame}

\begin{frame}{Studi osservazionali}
\phantomsection\label{studi-osservazionali}
Gli studi osservazionali danno informazioni importante, ma è sempre
dietro l'angolo il rischio di fare interpretazioni causali, \emph{specie
quando queste sono verosimili}.

Questo però non esclude la possibilità di spiegazioni alternative.
\end{frame}

\begin{frame}{Problemi studi osservazionali}
\phantomsection\label{problemi-studi-osservazionali}
Negli studi osservazionali c'è spesso la tendenza a pensare che la cosa
ideale sia ``controllare'' per covariate o inserire variabili nel
modello per isolare un predittore. Questo può però creare addirittura
errori, se non si consiedano certi effetti (es. \textbf{colliders})
\end{frame}

\begin{frame}{Sintesi di Evidenze}
\phantomsection\label{sintesi-di-evidenze}
Anche nella sintesi di evidenze attenzione a

``garbage in'' \(\to\) ``garbage out''
\end{frame}

\begin{frame}{\small Come può contribuire il personale clinico alla
ricerca? (2/2)}
\phantomsection\label{come-puuxf2-contribuire-il-personale-clinico-alla-ricerca-22}
\small

\begin{itemize}
\item
  Ricordando che per contribuire alla ricerca \textbf{non è necessario
  essere ricercatori a tempo pieno}: anche chi lavora prevalentemente in
  ambito clinico può avere un ruolo decisivo nel migliorare la qualità
  delle conoscenze disponibili.
\item
  Collaborando attivamente con i ricercatori, non solo come semplice
  ``fornitore di dati'', ma \textbf{contribuendo a definire e affinare
  la domanda di ricerca} a partire dai bisogni reali osservati nella
  pratica clinica.
\item
  \textbf{Aiutando a individuare quali problemi siano davvero rilevanti
  per i pazienti e per il lavoro quotidiano}, evitando studi che
  rispondano a domande poco utili o lontane dalla realtà clinica.
\item
  \textbf{Partecipando alla raccolta di dati di qualità}, accurati e
  standardizzati, elemento fondamentale per produrre risultati
  affidabili.
\end{itemize}
\end{frame}

\begin{frame}{\small Come può contribuire il personale clinico alla
ricerca? (2/2)}
\phantomsection\label{come-puuxf2-contribuire-il-personale-clinico-alla-ricerca-22-1}
\small

\begin{itemize}
\item
  Contribuendo all' \textbf{interpretazione dei risultati}, perché il
  punto di vista clinico è essenziale per capire se un risultato
  statisticamente significativo abbia anche un reale valore pratico.
\item
  Favorendo la collaborazione tra reparti, centri e professionisti
  diversi, così da costruire dataset più ampi e rappresentativi.
\item
  \textbf{Privilegiando progetti multicentrici, grandi database clinici}
  (possibilmente aperti e condivisibili) e trial ben disegnati,
  piuttosto che piccoli studi isolati che rischiano di produrre
  risultati fragili o poco utili.
\item
  \textbf{Promuovendo una cultura della valutazione critica delle
  evidenze}, aiutando colleghi e pazienti a interpretare i risultati
  della ricerca con prudenza.
\end{itemize}
\end{frame}




\end{document}
